\bookchapter{Capitolo 50: Kotlin a guida rapida — il primo linguaggio dello sviluppo Android}

\chapterintro{In questo capitolo iniziamo davvero con il linguaggio. Niente ansia — Kotlin è famoso per essere «amichevole con i principianti», e non devi imparare a memoria la grammatica: padroneggiando il 20\% della sintassi capisci e scrivi l'80\% del codice di ogni giorno. Il resto lo cercherai man mano che fai i progetti. Ricorda il nostro principio di studio: capire > memorizzare, saper usare > essere perfetti.}

\booksection{Variabili e tipi: dare un nome ai dati}

Un programma è «dati + logica». La variabile è il nome che diamo ai dati per conservarli. In Kotlin si dichiara in due modi:

\begin{lstlisting}[language=Kotlin]
val name = Marco        // val: variabile di sola lettura (il consigliato per default)
var age = 18             // var: variabile mutabile (si può cambiare dopo)

age = 19                 // <i class="fas fa-check-circle mr-1"></i> var si può cambiare
// name = "Anna"         // <i class="fas fa-times-circle mr-1"></i> val non si può cambiare: errore di compilazione

val score: Int = 100     // puoi anche dichiarare il tipo esplicitamente
val pi = 3.14            // se lo ometti viene dedotto automaticamente
\end{lstlisting}

I tipi più usati sono pochi: con una tabella li ricordi tutti:

\begin{center}
\begin{tabularx}{\linewidth}{|X|X|X|}
\hline
\textbf{Tipo} & \textbf{Cosa contiene} & \textbf{Esempio} \tabularnewline
\hline
Int & Numeri interi & 42 \tabularnewline
Double & Numeri decimali & 3.14 \tabularnewline
String & Testo & «Ciao» \tabularnewline
Boolean & Vero/falso & true / false \tabularnewline
List< T> & Un insieme di dati & listOf(1, 2, 3) \tabularnewline
\hline
\end{tabularx}
\end{center}

\begin{center}\small\color{gray}I tipi base più usati di Kotlin\end{center}

\begin{tcbinfo}{Consiglio: il principio «val prima di tutto»}
La regola d'oro della community Kotlin: per impostazione predefinita usa val, e var solo se serve davvero. Se usi var dove basterebbe val, l'IDE ti segnala in giallo. Questa abitudine ti risparmia tutta una famiglia di bug «variabile modificata per sbaglio» — meno cose cambiano nel codice, più è facile ragionarci.

\end{tcbinfo}

\booksection{Le funzioni: la logica impacchettata in «pulsanti»}

Una funzione è un pezzo di logica con un nome, come un pulsante del telecomando — lo premi (la chiami) e parte una serie di azioni. In Kotlin si dichiara con fun:

\begin{lstlisting}[language=Kotlin]
fun greet(name: String): String {
    return "Ciao, $name!"     // il simbolo $ inserisce la variabile nella stringa
}

println(greet(Marco))        // stampa: Ciao, Marco!

// le funzioni di una riga si possono scrivere come espressioni
fun add(a: Int, b: Int) = a + b

println(add(3, 5))            // stampa: 8
\end{lstlisting}

Qualche punto chiave: i parametri vanno tra parentesi (prima il nome, poi il tipo, separati dai due punti); il tipo di ritorno va dopo le parentesi; il template di stringa \$nomeVariabile è la comodità di punta di Kotlin — niente più incastri alla Java tra «stringa + variabile».

\booksection{Condizioni e cicli: far «pensare» il programma}

L'intelligenza di un programma nasce da due cose: saper decidere (if/when) e saper ripetere (for/while).

\begin{lstlisting}[language=Kotlin]
// Decisioni: if / else
val age = 17
if (age >= 18) {
    println("maggiorenne")
} else {
    println("minorenne")
}

// Multi-ramo: when (cento volte più comodo di switch)
val score = 85
val level = when {
    score >= 90 -> "ottimo"
    score >= 60 -> "sufficiente"
    else -> "Continua così!"
}

// Ripetizione: un ciclo for
val fruits = listOf("mela", "banana", "arancia")
for (fruit in fruits) {
    println("Mi piace la $fruit")
}
\end{lstlisting}

when è la costruzione a più rami più usata in Kotlin: nello sviluppo Android la userai ogni giorno (per esempio per mostrare interfacce diverse a seconda dello stato).

\booksection{La null safety: l'«amuleto salvavita» di Kotlin}

Il bug più celebre dell'era Java si chiama NullPointerException (eccezione di puntatore nullo) — il programma tocca qualcosa che «non c'è» e va subito in crash. Secondo le statistiche è da sempre primo nella classifica dei crash di Android. Kotlin lo risolve a livello di linguaggio: per impostazione predefinita i valori nulli non sono ammessi.

\begin{lstlisting}[language=Kotlin]
var name: String = "Marco"   // tipo normale: non può mai essere null
// name = null                // <i class="fas fa-times-circle mr-1"></i> errore di compilazione — ti blocca subito

var nickname: String? = null // aggiungi un punto di domanda: null consentito
nickname = "Nino"

// Con i tipi nullable devi gestirli in sicurezza:
println(nickname?.length)    // chiamata sicura: restituisce null invece di crashare
val len = nickname?.length ?: 0   // operatore Elvis: valore predefinito se null
println(len)
\end{lstlisting}

\begin{tcbwarn}{Attenzione: non confondere ? e ?}
String? significa «può essere nullo», a ?: b significa «se a è nullo usa b» (assomiglia alla pettinatura di Elvis, per questo si chiama operatore Elvis). Quando inizi a scrivere codice e il compilatore ti segnala «errori legati a null», nel 90\% dei casi hai dimenticato di gestire un tipo nullable — è proprio il linguaggio che ti sta proteggendo.

\end{tcbwarn}

\booksection{Classi e oggetti: gli stampi per fabbricare «pezzi»}

La classe è lo «stampo» degli oggetti, l'oggetto è il «prodotto finito» della classe. In Android tutto — schermate, pulsanti, dati — è un oggetto. In Kotlin c'è un'arma segreta chiamata data class (classe di dati), fatta apposta per «contenere dati»: una riga vale decine di righe di Java:

\begin{lstlisting}[language=Kotlin]
// Classe normale: uno stampo
class Dog(val name: String) {
    fun bark() = println("$name: Bau!")
}

val dog = Dog("Fido")      // crea un oggetto
dog.bark()                 // Fido: Bau!

// data class: una scatola per dati (con equals/copy incluse)
data class User(val id: Int, val name: String, val email: String)

val me = User(1, Marco, "zhangsan@example.com")
val twin = me.copy(name = "Marco Rossi")   // copia e cambia un campo
\end{lstlisting}

Quando farai il registro delle spese, ogni voce sarà una data class: data class Record(val amount: Double, val category: String, val date: String). Visto? Sei già capace di scrivere la struttura di codice di un progetto vero.

\booksection{Lambda e funzioni di ordine superiore: lo «zucchero sintattico» di Kotlin}

Nel codice Android vedrai ogni giorno «funzioni senza nome» come \{ ... \}: si chiamano lambda (funzioni anonime). onClick = \{ count++ \} su un Button è una lambda — impacchetta «cosa fare al tocco» in un blocco e lo consegna al pulsante.

\begin{lstlisting}[language=Kotlin]
// Modo classico: definisci una funzione con nome e passala
fun double(x: Int) = x * 2
val result = listOf(1, 2, 3).map(::double)

// Colle lambda: definita al volo, snella e diretta
val result2 = listOf(1, 2, 3).map { it * 2 }   // it = elemento corrente

// Combinazione filtro + trasformazione (frequentissima nei progetti reali)
val bills = listOf(12.5, 88.0, 5.0, 200.0)
val bigBills = bills.filter { it > 50 }.sorted()
// [88.0, 200.0] — "fatture oltre 50, ordinate"
\end{lstlisting}

Segna qualche «combo» ad alta frequenza sulle collezioni: filter (filtra), map (trasforma), sorted (ordina), sum (somma). Il 90\% dei lavori sulle liste che farai è una combinazione di queste quattro mosse. La filosofia di Kotlin è: se una logica si spiega in una frase, non scriverla mai in cinque.

\booksection{La scheda di sintassi rapida per chi parte da zero}

Non provare a imparare tutta la sintassi in un giorno. Salva questa scheda, consultala mentre scrivi codice e tra due settimane ti verrà naturale:

\begin{center}
\begin{tabularx}{\linewidth}{|X|X|}
\hline
\textbf{Cosa vuoi fare} & \textbf{Come si scrive} \tabularnewline
\hline
Definire una variabile di sola lettura & \texttt{val x = 1} \tabularnewline
Definire una variabile modificabile & \texttt{var x = 1} \tabularnewline
Definire una funzione & fun nome(parametri): TipoRitorno \{ ... \} \tabularnewline
Decidere & if (condizione) \{...\} else \{...\} \tabularnewline
Scelte multiple & when \{ condizione -> risultato \} \tabularnewline
Scorrere & for (x in lista) \{ ... \} \tabularnewline
Unire le stringhe & "Ciao, \$name" \tabularnewline
Gestire i valori nulli & \texttt{x?.\allowbreak length ?: 0} \tabularnewline
Definire una data class & \texttt{data class X(val a: Int, val b: String)} \tabularnewline
\hline
\end{tabularx}
\end{center}

\begin{center}\small\color{gray}Scheda di sintassi rapida Kotlin (versione principianti)\end{center}

\begin{tcbpractice}{Mani all'opera: gioca 10 minuti con Kotlin senza installare nulla}
Apri nel browser play.kotlinlang.org (il playground ufficiale online di Kotlin), incolla e fai girare ogni pezzo di codice di questo capitolo, cambia numeri e testi e guarda come cambia l'output. Questi 10 minuti di «prima si gioca, poi si studia» ti renderanno il capitolo 51 due volte più facile.

\end{tcbpractice}

\begin{tcbsummary}{Riepilogo del capitolo}
\begin{itemize}
  \item val prima di tutto, var con cautela; i tipi si deducono da soli, ma puoi dichiararli esplicitamente dopo i due punti
  \item Funzione = logica impacchettata; il template di stringa \$ è la comodità di punta
  \item if/when per decidere, for per ciclare; when è il «tuttofare» dello sviluppo Android
  \item La null safety è il salvavita di Kotlin: i tipi normali rifiutano null, quelli nullable vanno gestiti con ? / ?:
  \item La data class è una scatola per i dati, un campo per riga: nei progetti veri si usa ogni giorno
  \item Il 20\% della sintassi fa l'80\% del lavoro: prima impari a usarlo, poi diventi esperto
\end{itemize}

\end{tcbsummary}
